\section{Substituição de variáveis}
\label{sec:substituicao}
Para trabalharmos com quantificadores, precisaremos da noção de substituição de símbolos em fórmulas, feita como se segue.
Por razões técnicas e a fim de evitar fazer o mesmo trabalho duas vezes, também definiremos a substituição de constantes.

\begin{convention}\label{conv:mapaDeSubstituicao}
    Seja $\mathcal L$ um léxico para uma linguagem de primeira ordem.
    Um $\mathcal L$-mapeamento de variáveis e constantes é um mapeamento $\rho$ que associa cada elemento $x$ de alguma coleção finita $X$ de variáveis ou constantes de $\mathcal L$ a um $\mathcal L$-termo $\rho(x)$.
    A coleção $X$ é denominada \emph{domínio} de $\rho$, e dizemos que $\rho$ \emph{está definido} em $X$.

    Se $X$ possui apenas variáveis, dizemos que $\rho$ é um $\mathcal L$-mapeamento de variáveis.

    Se $X$ é, digamos, composto das variáveis ou constantes $x_1, \dots, x_n$, sem repetições, então $\rho$ pode ser denotada por $(x_1, \dots, x_n)\to (t_1, \dots, t_n)$, em que $t_i$ é $\rho(x_i)$ para cada $i\in\{1,\dots,n\}$.

    O mapeamento vazio é o mapeamento $\rho$ cujo domínio é a coleção vazia, ou seja, aquele que não associa nenhum elemento a nada.


    Dizemos que $\rho$ é um mapeamento \emph{trivial} se $\rho(x)$ é $x$ para cada $x \in X$.
\end{convention}

Com a convenção acima, podemos definir a substituição de variáveis e constantes em termos e fórmulas.
Iniciaremos com a substituição em termos.

\begin{metadefinition}[Substituição em termos]\label{mdef:substituicaoEmTermos}\index{substituição!em termos}
    Seja $\mathcal L$ um léxico para uma linguagem de primeira ordem, $X$ uma coleção finita de variáveis ou constantes, e $\rho$ um $\mathcal L$-mapeamento de variáveis e constantes definido em $X$.

    Dado um $\mathcal L$-termo $t$, $t[\rho]$ é o termo obtido a partir de $t$ substituindo-se todas as ocorrências de elementos $x$ de $X$ por $\rho(x)$.
    Formalmente, define-se $t[\rho]$ por indução sobre a estrutura de $t$, como se segue.
    \begin{itemize}
        \item Se $t$ é algum $x \in X$, então $t[\rho]$ é $\rho(x)$.
        \item Se $t$ é uma variável ou constante diferente de todo $x \in X$, então $t[\rho]$ é $t$.
        \item Se $t$ é $fe_1 \dots e_n$, em que $f$ é um símbolo funcional $n$-ário e $e_1, \dots, e_n$ são termos para os quais $e_1[\rho], \dots, e_n[\rho]$ já foram definidos, então $t[\rho]$ é $fe_1[\rho] \dots e_n[\rho]$.
    \end{itemize}

    Se $\rho$ é $(x_1, \dots, x_n)\to (t_1, \dots, t_n)$, a notação $t[x_1/t_1, \dots, x_n/t_n]$ denota $t[\rho]$.
\end{metadefinition}

Devemos provar que, na notação anterior, $t[\rho]$ é um $\mathcal L$-termo.

\begin{metalemma}\label{mlem:substituicaoEmTermos}
    Seja $\mathcal L$ um léxico para uma linguagem de primeira ordem, $t$ um $\mathcal L$-termo, e $\rho$ um $\mathcal L$-mapeamento de variáveis e constantes definido em $X$.
    Então:
    \begin{enumerate}[label=(\alph*)]
        \item $t[\rho]$ é um $\mathcal L$-termo.
        \item Para toda variável ou constante $z$, $z$ ocorre em $t[\rho]$ se, e somente se,
        \begin{itemize}
            \item $z$ ocorre em $t$ e não está em $X$, ou
            \item existe algum $x \in X$ tal que $z$ ocorre em $\rho(x)$ e $x$ ocorre em $t$.
        \end{itemize}
    \end{enumerate}
\end{metalemma}
\begin{proof}
    Provaremos simultaneamente os dois itens por indução sobre a estrutura de $t$.

    Suponha primeiro que $t$ é uma variável ou constante $w$.

    (a) Se $w\in X$, então $t[\rho]$ é $\rho(w)$, que é um $\mathcal L$-termo.
    Já se $w\notin X$, então $t[\rho]$ é $w$, que é um $\mathcal L$-termo.

    (b) Suponha que $z$ é uma variável ou constante que ocorre em $t[\rho]$.
    Nesse caso, $w=z$.
    Suponha, ainda que não existe algum $x\in X$ tal que $z$ ocorre em $\rho(x)$ e $x$ ocorre em $t$.
    Nesse caso, então $w=z$ não está em $X$, e $w=z$ ocorre em $t$.

    Reciprocamente, se $z$ é uma variável ou constante que ocorre em $t$ e não está em $X$, então $z$ é $w$ e $w$ não está em $X$, de modo que $t[\rho]$ é $w$, e, portanto, $z$ ocorre em $t[\rho]$.
    Já se existe algum $x\in X$ tal que $z$ ocorre em $\rho(x)$ e $x$ ocorre em $t$, então $x$ é $w$, e $t[\rho]$ é $\rho(w)=\rho(x)$, de modo que $z$ ocorre em $t[\rho]$.
    \smallskip
    Agora suponha que $t$ é $fe_1\dots e_n$, em que $f$ é um símbolo funcional $n$-ário com $n\geq 1$ e $e_1,\dots,e_n$ são termos.

    (a) Pela hipótese indutiva, $e_i[\rho]$ é um $\mathcal L$-termo para cada $i\in\{1,\dots,n\}$.
    Portanto,
    \[
        t[\rho]=fe_1[\rho]\dots e_n[\rho]
    \]
    é um $\mathcal L$-termo.

    (b) Uma variável ou constante $z$ ocorre em $t[\rho]$ se, e somente se, $z$ ocorre em algum $e_i[\rho]$.
    Pela hipótese indutiva, isso ocorre se, e somente se, para algum $i\in\{1,\dots,n\}$, uma das condições abaixo ocorre:
    \begin{itemize}
        \item $z$ ocorre em $e_i$ e $z\notin X$;
        \item existe algum $x\in X$ tal que $x$ ocorre em $e_i$ e $z$ ocorre em $\rho(x)$.
    \end{itemize}

    Como as ocorrências de variáveis e constantes em $t=fe_1\dots e_n$ são exatamente as ocorrências que aparecem em algum dos termos $e_i$, isso equivale a dizer que uma das condições abaixo ocorre:
    \begin{itemize}
        \item $z$ ocorre em $t$ e $z\notin X$, ou
        \item existe algum $x\in X$ tal que $x$ ocorre em $t$ e $z$ ocorre em $\rho(x)$.
    \end{itemize}

    Isso conclui a indução.
\end{proof}
Sigamos para alguns exemplos.

\begin{example}\label{exa:substituicaoEmTermos}
    Considere o léxico da teoria dos anéis do Exemplo~\ref{exa:exemplosLexicos}.
    \begin{itemize}
        \item $x_1[x_1/0]$ é $0$.
        \item $x[y/0]$ é $x$.
        \item $(x_1 + x_7)[x_1/0, x_7/x_2]$ é $0 + x_2$.
        \item $(x_1 + x_2)[x_1/x_2, x_2/(x_1 + x_2)]$ é $x_2 + (x_1 + x_2)$.
        \item $(x_1 + 0)[0/(x_1 + x_2)]$ é $x_1 + (x_1 + x_2)$.
    \end{itemize}
\end{example}

Convém estabelecer uma convenção mais usual sobre a notação de substituição em termos, compatível com a prática matemática usual.
\begin{convention}\label{conv:substituicaoEmTermos}
    Seja $\mathcal L$ um léxico para uma linguagem de primeira ordem, e $s=s(x_1, \dots, x_n)$ um $\mathcal L$-termo.

    Sejam $t_1, \dots, t_n$ $\mathcal L$-termos.
    
    Ao escrevermos $s(t_1, \dots, t_n)$, em que $t_1, \dots, t_n$ são $\mathcal L$-termos, estaremos nos referindo ao $\mathcal L$-termo $s[x_1/t_1, \dots, x_n/t_n]$.
\end{convention}

Refaçamos alguns dos exemplos anteriores e alguns novos.

\begin{example}\label{exa:substituicaoEmTermos2}
    Considere o léxico da teoria dos anéis do Exemplo~\ref{exa:exemplosLexicos}.
    Considere, em cada item, o termo $s$ indicado.
    \begin{itemize}
        \item Se $s(x_1)$ é $x_1$, então $s(0)$ é $0$.
        \item Se $s(x, y)$ é $x$, então $s(x, 0)$ é $x$.
        \item Se $s(x_1, x_7)$ é $x_1 + x_7$, então $s(0, x_2)$ é $0 + x_2$.
        \item Se $s(x_1, x_2)$ é $x_1 + x_2$, então $s(x_2, x_1+x_2)$ é $x_2 + (x_1 + x_2)$.
        \item Se $s(x, y)$ é $x+y$, então $s(z_1, z_2)$ é $z_1 + z_2$.
        \item Se $s(y, x)$ é $x+y$, então $s(z_1, z_2)$ é $z_2 + z_1$.
    \end{itemize}
\end{example}

É importante ressaltar que ao escrever que $s=s(x_1, \dots, x_n)$, estamos estabelecendo uma convenção de ordem para $x_1, \dots, x_n$, conforme ilustrado pelos dois últimos itens do exemplo acima.

A substituição em fórmulas traz um cuidado adicional: variáveis que aparecem nos termos inseridos podem passar a ficar no escopo de quantificadores já presentes na fórmula.

Optaremos por definir primeiro a substituição formal em fórmulas sem impedir esse fenômeno e, em seguida, definir a noção de \emph{substituição válida}, que delimita os casos em que tal captura não ocorre.

Primeiro, mais algumas convenções.
\begin{convention}\label{conv:substituicaoEmFormulas}
    Seja $\mathcal L$ um léxico para uma linguagem de primeira ordem, e seja $\rho$ um $\mathcal L$-mapeamento de variáveis e constantes definido em uma coleção finita $X$.

    Para cada variável $z$, escrevemos $X\setminus\{z\}$ para denotar a coleção dos elementos de $X$ distintos de $z$.

    Denotamos por $\rho_{-z}$ o $\mathcal L$-mapeamento definido em $X\setminus\{z\}$ que associa a cada elemento $x$ de $X\setminus\{z\}$ o termo $\rho(x)$.
\end{convention}

Na convenção acima, $z$ pode ou não estar em $X$.
\begin{metadefinition}[Substituição em fórmulas]\label{mdef:substituicaoEmFormulas}\index{substituição!em fórmulas}
    Seja $\mathcal L$ um léxico para uma linguagem de primeira ordem e $\rho$ um $\mathcal L$-mapeamento de variáveis e constantes definido em $X$.

    Dada uma fórmula $\varphi$, $\varphi[\rho]$ é a fórmula obtida por uma substituição literal: substituem-se as ocorrências livres das variáveis de $X$ e todas as ocorrências das constantes de $X$ por suas respectivas imagens por $\rho$.

    Formalmente, define-se $\varphi[\rho]$ por indução sobre a estrutura de $\varphi$ pelo algoritmo a seguir.
    Supondo definido $\psi[\rho_{-y}]$ e  $\psi[\rho]$ para toda subfórmula própria $\psi$ de $\varphi$ e toda variável $y$ sobre a qual age algum quantificador em $\varphi$, define-se $\varphi[\rho]$ como segue.
    \begin{itemize}
        \item Se $\varphi$ é $P e_1 \dots e_n$, em que $P$ é um símbolo de predicado $n$-ário ($n>0$) e $e_1, \dots, e_n$ são termos, então $\varphi[\rho]$ é $P e_1[\rho] \dots e_n[\rho]$.
        \item Se $\varphi$ é um símbolo de predicado 0-ário, então $\varphi[\rho]$ é $\varphi$.
        \item Se $\varphi$ é $\neg \psi$, então $\varphi[\rho]$ é $\neg \psi[\rho]$.
        \item Se $\varphi$ é $\psi \square \theta$, em que $\square$ é $\vee$, $\wedge$, $\rightarrow$ ou $\leftrightarrow$, então $\varphi[\rho]$ é $\psi[\rho] \square \theta[\rho]$.
        \item Se $\varphi$ é $\forall y \psi$ ou $\exists y \psi$, em que $y$ é uma variável, então $\varphi[\rho]$ é $\forall y \psi[\rho_{-y}]$ ou $\exists y \psi[\rho_{-y}]$.
    \end{itemize}

    Se $\rho$ é $(x_1, \dots, x_n)\to (t_1, \dots, t_n)$, a notação $\varphi[x_1/t_1, \dots, x_n/t_n]$ denota $\varphi[\rho]$.
\end{metadefinition}

Novamente, por uma indução, prova-se que se $\varphi$ é uma fórmula, então $\varphi[\rho]$ é uma fórmula.

\begin{metalemma}\label{mlem:substituicaoEmFormulas}
    Seja $\mathcal L$ um léxico para uma linguagem de primeira ordem, $\varphi$ uma $\mathcal L$-fórmula, e $\rho$ um $\mathcal L$-mapeamento de variáveis e constantes definido em $X$.
    Então $\varphi[\rho]$ é uma $\mathcal L$-fórmula.
\end{metalemma}
\begin{proof}

    Suponha primeiro que $\varphi$ é uma fórmula atômica $R(e_1, \dots, e_n)$, em que $R$ é um símbolo de predicado $n$-ário com $n\geq 1$ e $e_1, \dots, e_n$ são termos.

    Pelo Metalema~\ref{mlem:substituicaoEmTermos}, $e_i[\rho]$ é um $\mathcal L$-termo para cada $i\in\{1,\dots,n\}$.
    Segue que $\varphi[\rho]$ é $R(e_1[\rho], \dots, e_n[\rho])$, que é uma $\mathcal L$-fórmula.
    \smallskip

    Agora suponha que $\varphi$ é $R$, em que $R$ é um símbolo de predicado $0$-ário.
    
    $\varphi[\rho]$ é $R$, que é uma $\mathcal L$-fórmula.
    \smallskip

    Suponha que $\varphi$ é $\neg \psi$ e a tese já foi demonstrada para $\psi$.
    $\varphi[\rho]$ é $\neg \psi[\rho]$, que é uma $\mathcal L$-fórmula.
    \smallskip

    Suponha que $\varphi$ é $\psi \square \theta$, em que $\square$ é $\vee$, $\wedge$, $\rightarrow$ ou $\leftrightarrow$, e a tese já foi demonstrada para $\psi$ e para $\theta$.

    $\varphi[\rho]$ é $\psi[\rho] \square \theta[\rho]$, que é uma $\mathcal L$-fórmula.
    \smallskip

    Suponha que $\varphi$ é $Qy\psi$, em que $Q$ é $\forall$ ou $\exists$, e a tese já foi demonstrada para $\psi$.

    $\varphi[\rho]$ é $Qy\psi[\rho_{-y}]$, que é uma $\mathcal L$-fórmula.
\end{proof}

Note que a substituição ocorre mesmo que variáveis que ocorrem em $\rho(x)$ para algum $x$ em $X$ passem a ficar no escopo de quantificadores já presentes em $\varphi$.
Conforme veremos nos exemplos a seguir, tais substituições podem ser problemáticas.
A seguinte definição delimita as substituições em que tais problemas não ocorrem.

\begin{metadefinition}\label{mdef:substituibilidade}\index{substituição!válida}\index{substituível}
    Seja $\mathcal L$ um léxico para uma linguagem de primeira ordem, seja $\varphi$ uma $\mathcal L$-fórmula, seja $x$ uma variável ou constante e seja $t$ um $\mathcal L$-termo.
    Dizemos que $x$ é substituível por $t$ em  $\varphi$ se nenhuma ocorrência livre de $x$ em $\varphi$ estiver no escopo de um quantificador que age em alguma variável que ocorre em $t$.

    Se $\rho$ é um $\mathcal L$-mapeamento de variáveis e constantes definido em $X$, dizemos que \emph{$\rho$ é uma substituição válida em $\varphi$}, ou que \emph{a substituição $\varphi[\rho]$ é válida}, se cada $x \in X$ for substituível por $\rho(x)$ em $\varphi$.

    Lembrete: considera-se que toda ocorrência de uma constante é livre.
\end{metadefinition}

Ainda não apresentamos uma noção de verdade ou de dedução, mas, utilizando o conhecimento prévio do leitor, podemos justificar a Metadefinição~\ref{mdef:substituibilidade} com exemplos.

\begin{example}\label{exa:substituicaoInvalida}
    Considere a fórmula $\exists y \, \neg(x = y)$.

    A variável $x$ não é substituível pelo termo $y$ nesta fórmula, pois a ocorrência livre de $x$ está no escopo de um quantificador que age em uma variável que ocorre em $y$ (o próprio $y$).

    Notemos que, ao substituir $x$ por $y$ nessa fórmula, obtemos $\exists y \, \neg(y = y)$.
\end{example}

\begin{example}\label{exa:substituicaoSemVariaveisLivres}
    Considere a fórmula $\varphi$ dada por $\forall x \exists y \, \neg(x = y)$.

    A substituição de $x$ e $y$ por qualquer termo $t$ é válida, pois a fórmula não possui variáveis livres.
    Notemos que, para qualquer termo $t$, $\varphi[x/t]$ é $\forall x \exists y \, \neg(x = y)$ e $\varphi[y/t]$ é $\forall x \exists y \, \neg(x = y)$, ou seja, a fórmula é invariante por tais substituições.
\end{example}



O metalema abaixo nos mostra critérios úteis em induções para demonstrar que certas substituições são válidas.

\begin{metalemma}\label{mlem:criteriosSubstituibilidade}
    Seja $\mathcal L$ um léxico para uma linguagem de primeira ordem.
    Sejam $\varphi$ e $\psi$ $\mathcal L$-fórmulas e seja $\rho$ um $\mathcal L$-mapeamento de variáveis e constantes definido em $X$.
    Então:
    \begin{itemize}
        \item Se $\varphi$ é uma fórmula atômica, então a substituição $\varphi[\rho]$ é válida.
        \item A substituição $\varphi[\rho]$ é válida se, e somente se, a substituição $(\neg \varphi)[\rho]$ é válida.
        \item As substituições $\varphi[\rho]$ e $\psi[\rho]$ são válidas se, e somente se, a substituição $(\varphi\square\psi)[\rho]$ é válida, em que $\square$ é $\vee$, $\wedge$, $\rightarrow$ ou $\leftrightarrow$.
        \item Suponha que $z$ é uma variável, que $Q$ é $\forall$ ou $\exists$.
        Então a substituição $(Qz\varphi)[\rho]$ é válida se, e somente se, a substituição $\varphi[\rho_{-z}]$ é válida e, para todo $x\in X\setminus\{z\}$, se $x$ ocorre livre em $\varphi$, então $z$ não ocorre em $\rho(x)$.
    \end{itemize}
\end{metalemma}
\begin{proof}
    Se $\varphi$ é atômica, então $\varphi$ não possui quantificadores.
    Assim, a tese é imediata.

    Para a negação, seja $x \in X$ e suponha que $y$ é uma variável que ocorre em $\rho(x)$.
    Note que se $x$ ocorre livre no escopo de algum quantificador que age em $y$ em $\neg \varphi$, o mesmo ocorre em $\varphi$, e reciprocamente.
    Logo, a substituição $(\neg \varphi)[\rho]$ é válida se, e somente se, a substituição $\varphi[\rho]$ é válida.
    

    Para os conectivos binários, seja $x \in X$ e suponha que $y$ é uma variável que ocorre em $\rho(x)$.
    Note que $x$ ocorre livre no escopo de algum quantificador que age em $y$ em $\varphi \square \psi$ se, e somente se, isso ocorre em $\varphi$ ou em $\psi$, de modo que a substituição $(\varphi \square \psi)[\rho]$ é válida se, e somente se, as substituições $\varphi[\rho]$ e $\psi[\rho]$ são válidas.

    Resta o caso dos quantificadores.
    Suponha que $z$ é uma variável, que $Q$ é $\forall$ ou $\exists$, que $\varphi[\rho_{-z}]$ é válida e que, para todo $x\in X\setminus\{z\}$, se $x$ ocorre livre em $\varphi$, então $z$ não ocorre em $\rho(x)$.

    Seja $x \in X$ e suponha que $Q'y\chi$ é o escopo de um quantificador da fórmula $Qz\varphi$ em que $x$ ocorre livre com relação à fórmula $Qz\varphi$.
    Veremos que $y$ não ocorre em $\rho(x)$.

    Como $x$ ocorre livre em $Qz\varphi$, $x$ não é $z$.

    Se $Q'y\chi$ é a própria fórmula $Qz\varphi$, então $y$ é $z$.
    Além disso, $x\neq z$ ocorre livre em $\varphi$, e, por hipótese, $z$ não ocorre em $\rho(x)$.
    Logo, $y$ não ocorre em $\rho(x)$.

    Já se $Q'y\chi$ não é a própria fórmula $Qz\varphi$, então $Q'y\chi$ é o escopo de um quantificador da fórmula $\varphi$ em que $x\neq z$ ocorre livre com relação à fórmula $\varphi$.
    Como $\varphi[\rho_{-z}]$ é válida e $\rho_{-z}(x)=\rho(x)$, temos que $y$ não ocorre em $\rho(x)$.

    Reciprocamente, suponha que a substituição $(Qz\varphi)[\rho]$ é válida.
    Então, para todo $x\in X\setminus\{z\}$, se $x$ ocorre livre em $\varphi$, então $z$ não ocorre em $\rho(x)$.
    De fato, se $z$ ocorresse em $\rho(x)$, então uma ocorrência livre de $x$ em $\varphi$ seria uma ocorrência livre de $x$ em $Qz\varphi$ no escopo do quantificador inicial, que age em $z$.

    Vejamos agora que a substituição $\varphi[\rho_{-z}]$ é válida.
    Seja $x\in X\setminus\{z\}$ e suponha que $y$ é uma variável que ocorre em $\rho_{-z}(x)=\rho(x)$.

    Suponha, por absurdo, que $x$ ocorre livre no escopo de algum quantificador que age em $y$ em $\varphi$.
    Como $x\neq z$, uma tal ocorrência de $x$ também é livre em $Qz\varphi$.
    Assim, $x$ ocorre livre no escopo de um quantificador que age em uma variável que ocorre em $\rho(x)$ em $Qz\varphi$, contradizendo a validade de $(Qz\varphi)[\rho]$.
\end{proof}

\begin{metalemma}\label{mlem:substituicaoEmFormulas2}
    Seja $\mathcal L$ um léxico para uma linguagem de primeira ordem, $\varphi$ uma $\mathcal L$-fórmula, e $\rho$ um $\mathcal L$-mapeamento de variáveis e constantes definido em $X$.
    Suponha que a substituição $\varphi[\rho]$ é válida.
    Então para toda variável ou constante $z$, $z$ ocorre livre em $\varphi[\rho]$ se, e somente se,
    \begin{itemize}
        \item $z$ ocorre livre em $\varphi$ e não está em $X$, ou
        \item existe algum $x \in X$ tal que $z$ ocorre em $\rho(x)$ e $x$ ocorre livre em $\varphi$.
    \end{itemize}
\end{metalemma}
\begin{proof}
    Provaremos a afirmação por indução sobre a estrutura de $\varphi$.

    Suponha primeiro que $\varphi$ é uma fórmula atômica $R(e_1, \dots, e_n)$, em que $R$ é um símbolo de predicado $n$-ário com $n\geq 1$ e $e_1, \dots, e_n$ são termos.

    Uma variável ou constante $z$ ocorre livre em $\varphi[\rho]$ se, e somente se, $z$ ocorre em algum $e_i[\rho]$.
    Pelo Metalema~\ref{mlem:substituicaoEmTermos}, isso ocorre se, e somente se, para algum $i\in\{1,\dots,n\}$, uma das condições abaixo ocorre:
    \begin{itemize}
        \item $z$ ocorre em $e_i$ e $z\notin X$;
        \item existe algum $x\in X$ tal que $x$ ocorre em $e_i$ e $z$ ocorre em $\rho(x)$.
    \end{itemize}

    Como as ocorrências livres de variáveis e constantes em $R(e_1, \dots, e_n)$ são exatamente as ocorrências que aparecem em algum dos termos $e_i$, isso equivale à tese.
    \smallskip

    Agora suponha que $\varphi$ é $R$, em que $R$ é um símbolo de predicado $0$-ário.

    $R$ não possui variáveis ou constantes, de modo que $z$ não ocorre livre em $\varphi[\rho]$ para toda variável ou constante $z$.
    Além disso, nenhuma variável ou constante ocorre livre em $\varphi$.
    Logo, a tese é satisfeita.
    \smallskip

    Suponha que $\varphi$ é $\neg \psi$ e a tese já foi demonstrada para $\psi$.
    Como a substituição $(\neg\psi)[\rho]$ é válida, pelo Metalema~\ref{mlem:criteriosSubstituibilidade}, a substituição $\psi[\rho]$ é válida.

    Uma variável ou constante $z$ ocorre livre em $\varphi[\rho]$ se, e somente se, $z$ ocorre livre em $\psi[\rho]$.
    Pela hipótese indutiva, isso ocorre se, e somente se,
    \begin{itemize}
        \item $z$ ocorre livre em $\psi$ e $z\notin X$; ou
        \item existe algum $x\in X$ tal que $x$ ocorre livre em $\psi$ e $z$ ocorre em $\rho(x)$.
    \end{itemize}
    Como as ocorrências livres de variáveis e constantes em $\neg \psi$ correspondem às ocorrências livres de variáveis e constantes em $\psi$, segue a tese.
    \smallskip

    Suponha que $\varphi$ é $\psi \square \theta$, em que $\square$ é $\vee$, $\wedge$, $\rightarrow$ ou $\leftrightarrow$, e a tese já foi demonstrada para $\psi$ e para $\theta$.
    Como a substituição $(\psi\square\theta)[\rho]$ é válida, pelo Metalema~\ref{mlem:criteriosSubstituibilidade}, as substituições $\psi[\rho]$ e $\theta[\rho]$ são válidas.

    Uma variável ou constante $z$ ocorre livre em $\varphi[\rho]$ se, e somente se, $z$ ocorre livre em $\psi[\rho]$ ou $z$ ocorre livre em $\theta[\rho]$.
    Pela hipótese indutiva, isso ocorre se, e somente se, uma das condições abaixo ocorre:
    \begin{itemize}
        \item $z$ ocorre livre em $\psi$ e $z\notin X$;
        \item existe algum $x\in X$ tal que $x$ ocorre livre em $\psi$ e $z$ ocorre em $\rho(x)$;
        \item $z$ ocorre livre em $\theta$ e $z\notin X$;
        \item existe algum $x\in X$ tal que $x$ ocorre livre em $\theta$ e $z$ ocorre em $\rho(x)$.
    \end{itemize}
    Como as ocorrências livres de variáveis e constantes em $\psi \square \theta$ correspondem às ocorrências livres de variáveis e constantes em $\psi$ ou em $\theta$, segue a tese.
    \smallskip

    Suponha que $\varphi$ é $Qy\psi$, em que $Q$ é $\forall$ ou $\exists$, e a tese já foi demonstrada para $\psi$.
    Pela validade de $(Qy\psi)[\rho]$ e pelo Metalema~\ref{mlem:criteriosSubstituibilidade}, a substituição $\psi[\rho_{-y}]$ é válida e, para todo $x\in X\setminus\{y\}$, se $x$ ocorre livre em $\psi$, então $y$ não ocorre em $\rho(x)$.

    Se $z$ é a variável $y$, então $z$ não ocorre livre em $(Qy\psi)[\rho]$.
    Além disso, $y$ não ocorre livre em $Qy\psi$, e, se algum $x\in X$ ocorre livre em $Qy\psi$, então $x\in X\setminus\{y\}$ e $x$ ocorre livre em $\psi$, de modo que $y$ não ocorre em $\rho(x)$.
    Portanto, a tese vale nesse caso.

    Agora suponha que $z$ é diferente de $y$.
    Uma variável ou constante $z$ ocorre livre em $(Qy\psi)[\rho]$ se, e somente se, $z$ ocorre livre em $\psi[\rho_{-y}]$.
    Pela hipótese indutiva aplicada a $\psi$ e $\rho_{-y}$, isso ocorre se, e somente se,
    \begin{itemize}
        \item $z$ ocorre livre em $\psi$ e $z\notin X\setminus\{y\}$; ou
        \item existe algum $x\in X\setminus\{y\}$ tal que $x$ ocorre livre em $\psi$ e $z$ ocorre em $\rho(x)$.
    \end{itemize}
    Como $z$ é diferente de $y$, isso equivale a dizer que
    \begin{itemize}
        \item $z$ ocorre livre em $Qy\psi$ e $z\notin X$; ou
        \item existe algum $x\in X$ tal que $x$ ocorre livre em $Qy\psi$ e $z$ ocorre em $\rho(x)$.
    \end{itemize}
    Segue a tese.
\end{proof}

Agora vamos desenvolver a notação usada na prática.
Começaremos com alguns lemas simples.
\begin{metalemma}\label{mlem:substituicaoTrivialTermos}
    Seja $\mathcal L$ um léxico para uma linguagem de primeira ordem e seja $\rho$ um $\mathcal L$-mapeamento trivial de variáveis e constantes.
    Então, para todo $\mathcal L$-termo $t$, $t[\rho]$ é $t$.

    Em particular, $t[x/x]$ é $t$ para todo $\mathcal L$-termo $t$ e toda variável ou constante $x$, e, sendo $\varepsilon$ o mapeamento vazio, $t[\varepsilon]$ é $t$ para todo $\mathcal L$-termo $t$.
\end{metalemma}
\begin{proof}
    Prova-se por indução sobre a estrutura de $t$ que $t[\rho]=t$.
    Seja $X$ a coleção de variáveis ou constantes para a qual $\rho$ é definido.

    Se $z$ é uma variável ou constante, temos que $z[\rho]$ é $z$: se $z$ está em $X$, então $\rho(z)$ é $z$, e, se $z$ não está em $X$, então $z[\rho]$ é $z$ por definição.

    Se $t$ é $fe_1 \dots e_n$, em que $f$ é um símbolo funcional $n$-ário e $e_1, \dots, e_n$ são termos para os quais $e_1[\rho]=e_1$, \dots, $e_n[\rho]=e_n$, então $t[\rho]$ é $fe_1[\rho] \dots e_n[\rho]$, que é $fe_1 \dots e_n$.
\end{proof}

\begin{metalemma}\label{mlem:substituicaoTrivial}
    Seja $\mathcal L$ um léxico para uma linguagem de primeira ordem e seja $\rho$ um $\mathcal L$-mapeamento trivial de variáveis e constantes.
    Então, para toda $\mathcal L$-fórmula $\varphi$, $\varphi[\rho]$ é $\varphi$ e a substituição $\varphi[\rho]$ é válida.

    Em particular, $\varphi[x/x]$ é $\varphi$ para toda $\mathcal L$-fórmula $\varphi$ e toda variável ou constante $x$, e, sendo $\varepsilon$ o mapeamento vazio, $\varphi[\varepsilon]$ é $\varphi$ para toda $\mathcal L$-fórmula $\varphi$.
    Tais substituições são válidas.
\end{metalemma}

\begin{proof}
    Prova-se por indução sobre a estrutura de $\varphi$ que, para todo $\mathcal L$-mapeamento trivial $\rho$ de variáveis e constantes, $\varphi[\rho]$ é $\varphi$ e a substituição $\varphi[\rho]$ é válida.

    Considere uma fórmula atômica $R(e_1, \dots, e_n)$, em que $R$ é um símbolo de predicado $n$-ário com $n\geq 1$ e $e_1, \dots, e_n$ são termos.
    Por definição, $R(e_1, \dots, e_n)[\rho]$ é $R(e_1[\rho], \dots, e_n[\rho])$, que, pelo Metalema~\ref{mlem:substituicaoTrivialTermos}, é $R(e_1, \dots, e_n)$.
    A substituição $R(e_1, \dots, e_n)[\rho]$ é válida pelo Metalema~\ref{mlem:criteriosSubstituibilidade}.

    Se $R$ é um símbolo de predicado 0-ário, então $R[\rho]$ é $R$.
    A substituição $R[\rho]$ é válida pelo Metalema~\ref{mlem:criteriosSubstituibilidade}.

    Se $\varphi$ é $\neg \psi$ e a tese já foi demonstrada para $\psi$, então $\varphi[\rho]$ é $\neg \psi[\rho]$, que é $\neg \psi$.
    A substituição $(\neg \psi)[\rho]$ é válida pelo Metalema~\ref{mlem:criteriosSubstituibilidade}.

    Se $\varphi$ é $\psi \square \theta$, em que $\square$ é $\vee$, $\wedge$, $\rightarrow$ ou $\leftrightarrow$, e a tese já foi demonstrada para $\psi$ e para $\theta$, então $\varphi[\rho]$ é $\psi[\rho] \square \theta[\rho]$, que é $\psi \square \theta$.
    A substituição $(\psi\square\theta)[\rho]$ é válida pelo Metalema~\ref{mlem:criteriosSubstituibilidade}.
    
    Se $\varphi$ é $Q y \psi$, em que $Q$ é $\forall$ ou $\exists$ e $y$ é uma variável, e a tese já foi demonstrada para $\psi$, então $\varphi[\rho]$ é $Q y \psi[\rho_{-y}]$.
    Como $\rho_{-y}$ é um mapeamento trivial de variáveis e constantes definido em $X\setminus\{y\}$, em que $X$ é o domínio de $\rho$, segue da hipótese indutiva que $\psi[\rho_{-y}]$ é $\psi$.
    Logo, $\varphi[\rho]$ é $Q y \psi$, isto é, $\varphi$.
    A substituição $(Q y \psi)[\rho]$ é válida pelo Metalema~\ref{mlem:criteriosSubstituibilidade}, pois $\psi[\rho_{-y}]$ é válida por hipótese indutiva e, para todo $x\in X\setminus\{y\}$, se $x$ ocorre livre em $\psi$, então $y$ não ocorre em $\rho(x)$: de fato, $\rho(x)=x$, e $x$ é distinto de $y$ se for variável, enquanto, se $x$ for constante, o termo $x$ não possui variáveis.
\end{proof}

Extensões ``triviais'' de $\mathcal L$-mapeamentos não influenciam no resultado da substituição, como mostram os metalemas abaixo.

\begin{metalemma}\label{mlem:substituicaoExtensaoTrivialEmTermos}
    Seja $\mathcal L$ um léxico para uma linguagem de primeira ordem, $t$ um $\mathcal L$-termo e $\rho$ um $\mathcal L$-mapeamento de variáveis e constantes definido em $X$.

    Seja $\rho'$ um $\mathcal L$-mapeamento de variáveis e constantes definido em $X'$ estendendo $\rho$ tal que, para todo $x$ em $X'$ mas não em $X$:
    \begin{itemize}
        \item $x$ é uma variável ou constante que não ocorre em $t$, ou
        \item $\rho'(x)$ é $x$.
    \end{itemize}

    Então $t[\rho']$ é $t[\rho]$.
\end{metalemma}
\begin{proof}
    Prova-se por indução sobre a estrutura de $t$ que $t[\rho']=t[\rho]$.

    Se $t$ é uma variável ou constante $z$ em $X$, como $\rho'(z)=\rho(z)$, temos que $z[\rho']=\rho'(z)=\rho(z)=z[\rho]$.

    Se $t$ é uma variável ou constante $z$ em $X'$, mas não em $X$, então $z[\rho]$ é $z$ por definição, e $z[\rho']$ é $z$ por hipótese, já que, nesse caso, $\rho'(z)$ é $z$.

    Se $t$ é uma variável ou constante $z$ fora de $X'$, então $z[\rho]$ e $z[\rho']$ são $z$ por definição.

    Suponha agora que $t$ é $fe_1 \dots e_n$, em que $f$ é um símbolo funcional $n$-ário com $n\geq 1$ e $e_1, \dots, e_n$ são termos.
    Se a hipótese do enunciado é válida para $t$, então é válida para cada $e_i$.
    Logo, pela hipótese indutiva, $e_i[\rho']=e_i[\rho]$ para cada $i\in\{1,\dots,n\}$, e $t[\rho']$ é $fe_1[\rho'] \dots e_n[\rho']$, que é $fe_1[\rho] \dots e_n[\rho]$, isto é, $t[\rho]$.
\end{proof}

\begin{metalemma}\label{mlem:substituicaoTrivialEmFormulas}
    Seja $\mathcal L$ um léxico para uma linguagem de primeira ordem, $\varphi$ uma $\mathcal L$-fórmula e $\rho$ um $\mathcal L$-mapeamento de variáveis e constantes definido em $X$.

    Seja $\rho'$ um $\mathcal L$-mapeamento de variáveis e constantes definido em $X'$ estendendo $\rho$ tal que, para todo $x$ em $X'$ mas não em $X$:
    \begin{itemize}
        \item $x$ é uma variável ou constante que não ocorre livre em $\varphi$, ou
        \item $\rho'(x)$ é $x$.
    \end{itemize}

    Então $\varphi[\rho']$ é $\varphi[\rho]$ e a substituição $\varphi[\rho']$ é válida se, e somente se, a substituição $\varphi[\rho]$ é válida.
\end{metalemma}
\begin{proof}
    Prova-se por indução sobre a estrutura de $\varphi$ que $\varphi[\rho']=\varphi[\rho]$ e que a substituição $\varphi[\rho']$ é válida se, e somente se, a substituição $\varphi[\rho]$ é válida, quaisquer que sejam $\rho$ e $\rho'$ que satisfaçam as hipóteses do enunciado.

    Considere uma fórmula atômica $R(e_1,\dots,e_n)$, em que $R$ é um símbolo de predicado $n$-ário com $n\geq 1$ e $e_1,\dots,e_n$ são termos.
    Para cada $i\in\{1,\dots,n\}$, a hipótese do enunciado é válida para $e_i$.
    Logo, pelo Metalema~\ref{mlem:substituicaoExtensaoTrivialEmTermos}, temos $e_i[\rho']=e_i[\rho]$.
    Portanto, $R(e_1,\dots,e_n)[\rho']$ é $R(e_1[\rho'],\dots,e_n[\rho'])$, que é $R(e_1[\rho],\dots,e_n[\rho])$, isto é, $R(e_1,\dots,e_n)[\rho]$.
    Pelo Metalema~\ref{mlem:criteriosSubstituibilidade} ambas as substituições são válidas.

    Se $\varphi$ é um símbolo de predicado 0-ário, então $\varphi[\rho']$ e $\varphi[\rho]$ são ambas $\varphi$.
    Pelo Metalema~\ref{mlem:criteriosSubstituibilidade} ambas as substituições são válidas.

    Se $\varphi$ é $\neg\psi$, então a hipótese do enunciado é válida para $\psi$.
    Logo, pela hipótese indutiva, $\psi[\rho']=\psi[\rho]$.
    Portanto, $(\neg\psi)[\rho']$ é $\neg\psi[\rho']$, que é $\neg\psi[\rho]$, isto é, $(\neg\psi)[\rho]$.
    Também pela hipótese indutiva e pelo Metalema~\ref{mlem:criteriosSubstituibilidade}, a substituição $(\neg\psi)[\rho']$ é válida se, e somente se, a substituição $\psi[\rho']$ é válida, o que ocorre se, e somente se, a substituição $\psi[\rho]$ é válida, o que ocorre se, e somente se, a substituição $(\neg\psi)[\rho]$ é válida.

    Se $\varphi$ é $\psi\square\theta$, em que $\square$ é $\vee$, $\wedge$, $\rightarrow$ ou $\leftrightarrow$, então a hipótese do enunciado é válida para $\psi$ e para $\theta$.
    Logo, pela hipótese indutiva, $\psi[\rho']=\psi[\rho]$ e $\theta[\rho']=\theta[\rho]$.
    Portanto, $(\psi\square\theta)[\rho']$ é $\psi[\rho']\square\theta[\rho']$, que é $\psi[\rho]\square\theta[\rho]$, isto é, $(\psi\square\theta)[\rho]$.
    Também pela hipótese indutiva e pelo Metalema~\ref{mlem:criteriosSubstituibilidade}, a substituição $(\psi\square\theta)[\rho']$ é válida se, e somente se, as substituições $\psi[\rho']$ e $\theta[\rho']$ são válidas, o que ocorre se, e somente se, as substituições $\psi[\rho]$ e $\theta[\rho]$ são válidas, o que ocorre se, e somente se, a substituição $(\psi\square\theta)[\rho]$ é válida.

    Resta o caso em que $\varphi$ é $Qz\psi$, em que $Q$ é $\forall$ ou $\exists$ e $z$ é uma variável.
    Note que $\rho'_{-z}$ estende $\rho_{-z}$.
    Além disso, a hipótese do enunciado é válida para $\psi$, $\rho_{-z}$ e $\rho'_{-z}$.
    Logo, pela hipótese indutiva, $\psi[\rho'_{-z}]$ é $\psi[\rho_{-z}]$.
    Portanto, $(Qz\psi)[\rho']$ é $Qz\psi[\rho'_{-z}]$, que é $Qz\psi[\rho_{-z}]$, isto é, $(Qz\psi)[\rho]$.
    Também pela hipótese indutiva e pelo Metalema~\ref{mlem:criteriosSubstituibilidade}, a substituição $(Qz\psi)[\rho']$ é válida se, e somente se, a substituição $\psi[\rho'_{-z}]$ é válida e, para todo $x\in X'\setminus\{z\}$, se $x$ ocorre livre em $\psi$, então $z$ não ocorre em $\rho'(x)$.

    Como $\psi[\rho'_{-z}]$ é válida se, e somente se, $\psi[\rho_{-z}]$ é válida, resta comparar as condições sobre os elementos de $X'\setminus\{z\}$ e de $X\setminus\{z\}$.

    Para $x\in X\setminus\{z\}$, temos $\rho'(x)=\rho(x)$.
    Por outro lado, se $x$ está em $X'$ mas não em $X$, $x\neq z$ e $x$ ocorre livre em $\psi$, então $x$ ocorre livre em $Qz\psi$.
    Logo, pela hipótese sobre $\rho'$, temos $\rho'(x)=x$.
    Como $x$ é distinto de $z$, segue que $z$ não ocorre em $\rho'(x)$.

    Portanto, a condição
    \[
        \text{para todo }x\in X'\setminus\{z\},\text{ se }x\text{ ocorre livre em }\psi,\text{ então }z\text{ não ocorre em }\rho'(x)
    \]
    é equivalente à condição
    \[
        \text{para todo }x\in X\setminus\{z\},\text{ se }x\text{ ocorre livre em }\psi,\text{ então }z\text{ não ocorre em }\rho(x).
    \]

    Pelo Metalema~\ref{mlem:criteriosSubstituibilidade}, isso ocorre se, e somente se, a substituição $(Qz\psi)[\rho]$ é válida.
\end{proof}

\begin{convention}
    Seja $\mathcal L$ um léxico para uma linguagem de primeira ordem, $\rho$ um $\mathcal L$-mapeamento de variáveis e constantes definido em $X$ e $\sigma$ um $\mathcal L$-mapeamento de variáveis e constantes.

    O mapeamento $\rho\sigma$ é o $\mathcal L$-mapeamento de variáveis e constantes definido em $X$ tal que, para todo $x\in X$, $(\rho\sigma)(x)$ é $\rho(x)[\sigma]$.
\end{convention}

Nesse contexto, vale o seguinte metalema.

\begin{metalemma}\label{mlem:substituicaoCompostaTermos}
    Seja $\mathcal L$ um léxico para uma linguagem de primeira ordem e $t$ um $\mathcal L$-termo.
    Seja $\rho$ um $\mathcal L$-mapeamento de variáveis e constantes definido em $X$ e seja $\sigma$ um $\mathcal L$-mapeamento de variáveis e constantes definido em $Y$.

    Suponha que nenhum elemento de $Y$ ocorre em $t$.
    Então $t[\rho\sigma]$ é $t[\rho][\sigma]$.
\end{metalemma}

\begin{proof}
    Provaremos por indução sobre a estrutura de $t$.

    Se $t$ é uma variável ou constante $z$ fora de $X$, temos que $z[\rho\sigma]$ é $z$ por definição, e $z[\rho][\sigma]$ é $z[\sigma]$, que é $z$ por hipótese, já que nenhum elemento de $Y$ ocorre em $t$.

    Se $t$ é uma variável ou constante $z$ em $X$, então $z[\rho\sigma]$ é $\rho(z)[\sigma]$, que é $(\rho\sigma)(z)$ por definição.
    Portanto, $z[\rho][\sigma]$ é $z[\rho\sigma]$.

    Se $t$ é $fe_1\dots e_n$ e nenhum elemento de $Y$ ocorre em $t$, então nenhum elemento de $Y$ ocorre em cada $e_i$.
     Logo, pela hipótese indutiva aplicada a $e_1,\dots,e_n$, $e_i[\rho\sigma]$ é $e_i[\rho][\sigma]$ para cada $i\leq n$.
     Portanto,
    \[
        t[\rho\sigma]
        =
        fe_1[\rho\sigma]\dots e_n[\rho\sigma]
        =
        fe_1[\rho][\sigma]\dots e_n[\rho][\sigma]
        =(fe_1[\rho]\dots e_n[\rho])[\sigma]
        =(fe_1\dots e_n)[\rho][\sigma]
        =t[\rho][\sigma].
    \]
\end{proof}

\begin{metalemma}\label{mlem:substituicaoComposta}
    Seja $\mathcal L$ um léxico para uma linguagem de primeira ordem e $\varphi$ uma $\mathcal L$-fórmula.
    Seja $\rho$ um $\mathcal L$-mapeamento de variáveis e constantes definido em $X$ e seja $\sigma$ um $\mathcal L$-mapeamento de variáveis e constantes definido em $Y$.

    Suponha que nenhum elemento de $Y$ ocorre livre em $\varphi$ e que a substituição $\varphi[\rho]$ é válida.
    Então $\varphi[\rho\sigma]$ é $\varphi[\rho][\sigma]$.

    Além disso, a substituição $\varphi[\rho][\sigma]$ é válida se, e somente se, a substituição $\varphi[\rho\sigma]$ é válida.
\end{metalemma}

\begin{proof}
    Provaremos a afirmação por indução sobre a estrutura de $\varphi$.

    Se $\varphi$ é $R e_1\dots e_n$, em que $R$ é um símbolo de predicado $n$-ário ($n\geq 1$) e $e_1,\dots,e_n$ são termos, então, do Metalema~\ref{mlem:substituicaoCompostaTermos}, $e_i[\rho\sigma]$ é $e_i[\rho][\sigma]$ para cada $i\in\{1,\dots,n\}$.
    Logo,
    \[
        \varphi[\rho\sigma]
        =
        R e_1[\rho\sigma]\dots e_n[\rho\sigma]
        =
        R e_1[\rho][\sigma]\dots e_n[\rho][\sigma]
        =(R e_1[\rho]\dots e_n[\rho])[\sigma]
        =(R e_1\dots e_n)[\rho][\sigma]
        =\varphi[\rho][\sigma].
    \]
    Pelo Metalema~\ref{mlem:criteriosSubstituibilidade}, ambas as substituições são válidas.

    Se $\varphi$ é $R$, em que $R$ é um símbolo de predicado $0$-ário, então ambas as substituições $\varphi[\rho\sigma]$ e $\varphi[\rho][\sigma]$ são $R$.
    Pelo Metalema~\ref{mlem:criteriosSubstituibilidade}, ambas as substituições são válidas.


    Se $\varphi$ é $\neg\psi$, então, pela hipótese indutiva,
    \[
        \varphi[\rho\sigma]
        =
        \neg\psi[\rho\sigma]
        =
        \neg\psi[\rho][\sigma]
        =
        (\neg\psi)[\rho][\sigma]
        =
        \varphi[\rho][\sigma].
    \]
    Pelo Metalema~\ref{mlem:criteriosSubstituibilidade}, a substituição $\varphi[\rho\sigma]$ é válida se, e somente se, a substituição $\psi[\rho\sigma]$ é válida.
    Também pelo Metalema~\ref{mlem:criteriosSubstituibilidade}, a substituição $\varphi[\rho][\sigma]$ é válida se, e somente se, a substituição $\psi[\rho][\sigma]$ é válida.
    Como $\varphi[\rho]$ é válida, a substituição $\psi[\rho]$ é válida.
    Além disso, nenhum elemento de $Y$ ocorre livre em $\psi$.
    Assim, pela hipótese indutiva, a substituição $\psi[\rho\sigma]$ é válida se, e somente se, a substituição $\psi[\rho][\sigma]$ é válida.
    Portanto, a substituição $\varphi[\rho\sigma]$ é válida se, e somente se, a substituição $\varphi[\rho][\sigma]$ é válida.

    Se $\varphi$ é $\psi\square\theta$, em que $\square$ é $\vee$, $\wedge$, $\rightarrow$ ou $\leftrightarrow$, então, pela hipótese indutiva,
    \[
        \varphi[\rho\sigma]
        =
        \psi[\rho\sigma]\square\theta[\rho\sigma]
        =
        \psi[\rho][\sigma]\square\theta[\rho][\sigma]
        =
        (\psi\square\theta)[\rho][\sigma]
        =
        \varphi[\rho][\sigma].
    \]
    Pelo Metalema~\ref{mlem:criteriosSubstituibilidade}, a substituição $\varphi[\rho\sigma]$ é válida se, e somente se, as substituições $\psi[\rho\sigma]$ e $\theta[\rho\sigma]$ são válidas.
    Também pelo Metalema~\ref{mlem:criteriosSubstituibilidade}, a substituição $\varphi[\rho][\sigma]$ é válida se, e somente se, as substituições $\psi[\rho][\sigma]$ e $\theta[\rho][\sigma]$ são válidas.
    Como $\varphi[\rho]$ é válida, as substituições $\psi[\rho]$ e $\theta[\rho]$ são válidas.
    Além disso, nenhum elemento de $Y$ ocorre livre em $\psi$ ou em $\theta$.
    Assim, pela hipótese indutiva, $\psi[\rho\sigma]$ é válida se, e somente se, $\psi[\rho][\sigma]$ é válida, e $\theta[\rho\sigma]$ é válida se, e somente se, $\theta[\rho][\sigma]$ é válida.
    Portanto, a substituição $\varphi[\rho\sigma]$ é válida se, e somente se, a substituição $\varphi[\rho][\sigma]$ é válida.

    Resta o caso em que $\varphi$ é $Qz\psi$, em que $Q$ é $\forall$ ou $\exists$ e $z$ é uma variável.

    Como $(Qz\psi)[\rho]$ é válida, pelo Metalema~\ref{mlem:criteriosSubstituibilidade}, temos que $\psi[\rho_{-z}]$ é válida e, para todo $x\in X\setminus\{z\}$, se $x$ ocorre livre em $\psi$, então $z$ não ocorre em $\rho(x)$.

    Como nenhum elemento de $Y$ ocorre livre em $Qz\psi$, nenhum elemento de $Y\setminus\{z\}$ ocorre livre em $\psi$.
    Pela hipótese indutiva aplicada a $\psi$, $\rho_{-z}$ e $\sigma_{-z}$, temos
    \[
        \psi[\rho_{-z}\sigma_{-z}]
        =
        \psi[\rho_{-z}][\sigma_{-z}].
    \]

    Seja $X_0$ a coleção dos elementos de $X\setminus\{z\}$ que ocorrem livres em $\psi$.
    Para todo $x\in X_0$, temos que $z$ não ocorre em $\rho(x)$.
    Logo, pelo Metalema~\ref{mlem:substituicaoExtensaoTrivialEmTermos}, temos
    \[
        \rho(x)[\sigma]=\rho(x)[\sigma_{-z}].
    \]

    Seja $\tau$ o $\mathcal L$-mapeamento de variáveis e constantes definido em $X_0$ tal que, para cada $x\in X_0$,
    \[
        \tau(x)=\rho(x)[\sigma]=\rho(x)[\sigma_{-z}].
    \]
    Então $(\rho\sigma)_{-z}$ e $\rho_{-z}\sigma_{-z}$ são extensões de $\tau$ por elementos que não ocorrem livres em $\psi$.
    Portanto, pelo Metalema~\ref{mlem:substituicaoTrivialEmFormulas}, temos:
    \[
        \psi[(\rho\sigma)_{-z}]
        =
        \psi[\rho_{-z}\sigma_{-z}].
    \]

    Consequentemente,
    \[
        (Qz\psi)[\rho\sigma]
        =
        Qz\psi[(\rho\sigma)_{-z}]
        =
        Qz\psi[\rho_{-z}\sigma_{-z}]
        =
        Qz\psi[\rho_{-z}][\sigma_{-z}]
        =
        (Qz\psi)[\rho][\sigma].
    \]

    Resta comparar as condições de validade.

    Pelo Metalema~\ref{mlem:criteriosSubstituibilidade}, a substituição $(Qz\psi)[\rho\sigma]$ é válida se, e somente se:
    \begin{enumerate}[label=(\alph*)]
        \item a substituição $\psi[(\rho\sigma)_{-z}]$ é válida;
        \item para todo $x\in X\setminus\{z\}$, se $x$ ocorre livre em $\psi$, então $z$ não ocorre em $(\rho\sigma)(x)=\rho(x)[\sigma]$.
    \end{enumerate}

    Também pelo Metalema~\ref{mlem:criteriosSubstituibilidade}, a substituição $(Qz\psi)[\rho][\sigma]$ é válida se, e somente se:
    \begin{enumerate}[label=(\alph*')]
        \item a substituição $\psi[\rho_{-z}][\sigma_{-z}]$ é válida;
        \item para todo $u\in Y\setminus\{z\}$, se $u$ ocorre livre em $\psi[\rho_{-z}]$, então $z$ não ocorre em $\sigma(u)$.
    \end{enumerate}

    Pela hipótese indutiva aplicada a $\psi$, $\rho_{-z}$ e $\sigma_{-z}$, a substituição $\psi[\rho_{-z}\sigma_{-z}]$ é válida se, e somente se, a substituição $\psi[\rho_{-z}][\sigma_{-z}]$ é válida.
    Como $(\rho\sigma)_{-z}$ e $\rho_{-z}\sigma_{-z}$ são extensões de $\tau$ por elementos que não ocorrem livres em $\psi$, segue do Metalema~\ref{mlem:substituicaoTrivialEmFormulas} que a substituição $\psi[(\rho\sigma)_{-z}]$ é válida se, e somente se, a substituição $\psi[\rho_{-z}\sigma_{-z}]$ é válida.
    Assim, ($a$) é equivalente a ($a'$).

    Portanto, resta apenas comparar as condições ($b$) e ($b'$).

    Suponha primeiro ($b'$).
    Seja $x\in X\setminus\{z\}$ tal que $x$ ocorre livre em $\psi$.
    Como $(Qz\psi)[\rho]$ é válida, temos que $z$ não ocorre em $\rho(x)$.
    Suponha, por absurdo, que $z$ ocorre em $\rho(x)[\sigma]$.
    Pelo Metalema~\ref{mlem:substituicaoEmTermos}, existe algum $u\in Y$ tal que $u$ ocorre em $\rho(x)$ e $z$ ocorre em $\sigma(u)$.
    Como $z$ não ocorre em $\rho(x)$, temos $u\neq z$.
    Logo, $u\in Y\setminus\{z\}$.
    Como $\psi[\rho_{-z}]$ é válida, segue do Metalema~\ref{mlem:substituicaoEmFormulas2} que $u$ ocorre livre em $\psi[\rho_{-z}]$.
    Assim, $z$ não ocorre em $\sigma(u)$, contradição.
    Portanto, $z$ não ocorre em $\rho(x)[\sigma]$.

    Reciprocamente, suponha ($b$).
    Seja $u\in Y\setminus\{z\}$ tal que $u$ ocorre livre em $\psi[\rho_{-z}]$.
    Como $\psi[\rho_{-z}]$ é válida, pelo Metalema~\ref{mlem:substituicaoEmFormulas2}, uma das duas possibilidades ocorre:
    \begin{itemize}
        \item $u$ ocorre livre em $\psi$ e $u\notin X\setminus\{z\}$;
        \item existe algum $x\in X\setminus\{z\}$ tal que $x$ ocorre livre em $\psi$ e $u$ ocorre em $\rho(x)$.
    \end{itemize}

    A primeira possibilidade é impossível, pois nenhum elemento de $Y\setminus\{z\}$ ocorre livre em $\psi$.
    Logo, existe algum $x\in X\setminus\{z\}$ tal que $x$ ocorre livre em $\psi$ e $u$ ocorre em $\rho(x)$.

    Se $z$ ocorresse em $\sigma(u)$, então, pelo Metalema~\ref{mlem:substituicaoEmTermos}, $z$ ocorreria em $\rho(x)[\sigma]$, contrariando a hipótese.
    Portanto, $z$ não ocorre em $\sigma(u)$.
    Desse modo, a condição ($b'$) é satisfeita.
\end{proof}


\begin{metalemma}\label{mlem:substituicaoInversaEmFormulas}
    Seja $\mathcal L$ um léxico para uma linguagem de primeira ordem.

    Seja $\rho$ um $\mathcal L$-mapeamento de variáveis e constantes definido em $X$ e seja $\sigma$ um $\mathcal L$-mapeamento de variáveis e constantes definido em $Y$.
    Suponha que $\rho\sigma$ é trivial, isto é, que, para todo $x\in X$,
    \[
        \rho(x)[\sigma]=x.
    \]

    \begin{enumerate}
        \item Seja $t$ um $\mathcal L$-termo tal que nenhum elemento de $Y$ ocorre em $t$.
        Então $t[\rho][\sigma]$ é $t$.

        \item Seja $\varphi$ uma $\mathcal L$-fórmula tal que nenhum elemento de $Y$ ocorre livre em $\varphi$.
        Se a substituição $\varphi[\rho]$ é válida, então $\varphi[\rho][\sigma]$ é $\varphi$ e a substituição $\varphi[\rho][\sigma]$ é válida.
    \end{enumerate}
\end{metalemma}
\begin{proof}
    (a) Pelo Metalema~\ref{mlem:substituicaoCompostaTermos}, $t[\rho][\sigma]$ é $t[\rho\sigma]$.
    Como $\rho\sigma$ é trivial, $t[\rho\sigma]$ é $t$ pelo Metalema~\ref{mlem:substituicaoTrivialTermos}.
    Logo, $t[\rho][\sigma]$ é $t$.

    (b) Pelo Metalema~\ref{mlem:substituicaoComposta}, temos que $\varphi[\rho][\sigma]$ é $\varphi[\rho\sigma]$.
    Como $\rho\sigma$ é trivial, $\varphi[\rho\sigma]$ é $\varphi$ pelo Metalema~\ref{mlem:substituicaoTrivial}.
    Logo, $\varphi[\rho][\sigma]$ é $\varphi$.

    Além disso, pelo Metalema~\ref{mlem:substituicaoComposta}, a substituição $\varphi[\rho][\sigma]$ é válida se, e somente se, a substituição $\varphi[\rho\sigma]$ é válida.
    Como $\rho\sigma$ é trivial, a substituição $\varphi[\rho\sigma]$ é válida pelo Metalema~\ref{mlem:substituicaoTrivial}.
    Portanto, a substituição $\varphi[\rho][\sigma]$ é válida.
\end{proof}
\subsection*{Exercícios}

\begin{exercise}\label{exr:substituicaoEmFormulas}
    Considere a linguagem da teoria dos anéis do Exemplo~\ref{exa:exemplosLexicos}.
    Em cada item, determine $\varphi[x/t]$ e se a substituição é válida.
    \begin{enumerate}[label=(\alph*)]
        \item $\varphi$ é $x = x \cdot y$ e $t$ é $y$.
        \item $\varphi$ é $\forall x \, (x = x \cdot y)$ e $t$ é $y$.
        \item $\varphi$ é $\forall x \, (x = x \cdot y \wedge \forall y\, \neg (x = x))$ e $t$ é $y$.
        \item $\varphi$ é $ \neg (x+y =y+x) \wedge \forall y \, (x = x \cdot y)$ e $t$ é $y$.
    \end{enumerate}
\end{exercise}
